\section{No reflection-containing inducing block in six quadratic fields}
\label{sec:large-prime-inducing}

\begin{theorem}[D-BI-13: a ramification obstruction for every inducing block]
\label{D-BI-13}
Let
\[
 p\in\{7,11,19,43,67,163\},\qquad K=\mathbf Q(\sqrt{-p}).
\]
Suppose a finite-cover Weil sixfold source has exact character field
$K$, period rank nine, and nonsplit actual polarization.  If its
complex coefficient representation is expressed as
\[
                       \rho=\operatorname{Ind}_H^G\tau,
\]
then the finite image of $\tau$ contains no complex reflection.
The inducing representation may have arbitrary degree and arbitrary
character field.  In particular the obstruction applies to every
inducing block, not only to the original $K$-representation or to
blocks already defined over $K$.
\end{theorem}
\begin{proof}
Suppose the image of $\tau$ contains a reflection.  Let $R$ be
the normal subgroup generated by its reflections.  Its fixed space
is stable under the irreducible inducing image and hence is zero.
The nonorthogonality components of the reflection root lines give
a direct sum of irreducible reflection spaces permuted transitively
by that image.  Stabilize one component, of dimension $e$, and
denote its stabilizer representation by $\sigma$.  Transitivity
of induction gives
\[
                    \rho=\operatorname{Ind}_{H'}^G\sigma
\]
for a subgroup $H'\le H$.  Its faithful image $J$ contains an
irreducible reflection group $R_0$ of rank $e$ and normalizes it.
No rationality of this component or its field is asserted.  The
induced-character formula nevertheless gives the unconditional
inclusion
\[
                    K\subset L=\mathbf Q(\chi_\sigma).
                                                               \tag{BI.39}
\]

Ranks one, two, three and four are excluded, respectively, by
the monomial theorem D-MON-01--03, D-BI-01, D-BI-05, and
D-BI-11.  For $p=7$, the rank-three step uses the proved split
actual polarization of the $G_{24}$ family, with the polarized
Shapiro comparison in D-BI-05; this is where the nonsplit
hypothesis remains essential.  These are all-induction results with the larger field
retained, so they apply to $\sigma$.
If $R_0$ is imprimitive of rank $e\ge5$, Cohen's normalizer
Proposition~2.13 puts $J$ in the scalar monomial group
\cite[pp.~391--392]{FIELD:Cohen}; its low-rank exceptions have
already been discarded.  Then $\sigma$, and hence $\rho$, is
monomial, another contradiction.  If $R_0$ is a real group of
type $A_e$, its scalar normalizer is the deleted permutation
normalizer and D-BI-08 applies.  This also covers every rank in
that infinite family, without assuming that the scalar and
permutation parts form a direct product.

The published reflection classification now leaves precisely
\[
\begin{array}{c|c|c|c}
 R_0&e&\text{finite normalizer model }N\text{ up to scalars}&F\\\hline
 G_{33}&5&G_{33}&\mathbf Q(\zeta_3)\\
 G_{34}&6&G_{34}&\mathbf Q(\zeta_3)\\
 W(E_6)&6&\langle W(E_6),-I\rangle&\mathbf Q\\
 W(E_7)&7&W(E_7)&\mathbf Q\\
 W(E_8)&8&W(E_8)&\mathbf Q.
\end{array}                                                    \tag{BI.40}
\]
The classification is Cohen's Theorem~5.12, p.~435, and the
model fields are in Malle's Table~8.2, p.~675, with its explanation
on p.~676.  The omitted $G_{35}$ is the rational Weyl group
$W(E_6)$.  For $G_{33},G_{34}$ the normalizer assertion is
Cohen's corrected Proposition~5.14; its sole exception is rank
three and is already irrelevant
\cite{FIELD:Cohen,FIELD:CohenErr,FIELD:Malle}.
For the real rows, a normalizer is scalar times real by
Cohen's Lemma~2.12.  A root-system automorphism is a Weyl element
times a diagram automorphism.  For $E_6$ the diagram involution
is represented by $-w_0$, so adjoining $-I$ suffices; $E_7,E_8$
have no nontrivial diagram automorphisms and already contain
$-I$.  These are finite rational matrix groups.  Thus in every
row
\[
                    J\subset\mathbf C^*N,
                    \qquad N\subset\operatorname{GL}_e(F),
                    \qquad e\le8.                          \tag{BI.41}
\]

We next verify the branch count on which the argument depends.
Use the actual intermediate connected cover for $\sigma$.
Its cohomological multiplicity is six by Frobenius reciprocity,
and its full quadratic-field period rank factors through that
cover's marked-base moduli.  If the base has genus $h$, then
$D=6-2e(h-1)$.  Genus one gives at most six parameters, and
genus at least two is impossible because $e\ge5$.  Hence rank
nine forces a rational base and a product-one generating tuple
$g_1,\ldots,g_r$ of nonidentity inertia matrices with
\[
                  r\ge12,
                  \qquad\sum_j\operatorname{rank}(1-g_j)=2e+6.
                                                               \tag{BI.42}
\]

Write $g_j=\lambda_j n_j$ with $n_j\in N$.  Since both $g_j$
and $n_j$ have finite order and the scalar commutes with $n_j$,
each $\lambda_j$ is a root of unity.  The ambiguity in its choice
is exactly a scalar in $Z_0=N\cap\mathbf C^*I$.  Such a scalar
belongs to $F$, so its order divides two or six in (BI.40), and
in particular is prime to $p$.
It follows that the $p$-primary part of $\lambda_j$ is independent
of the chosen expression.  More generally these primary parts
define a character
\[
              \beta_p:J\longrightarrow\mu_{p^a}
                                                               \tag{BI.43}
\]
for some finite $a\ge0$.  One can see the homomorphism directly
from multiplication, or from the scalar quotient
$\mathbf C^*N/N\simeq\mathbf C^*/Z_0$ and its finite cyclic
image on $J$.  No scalar projection on $J$ before quotienting
$Z_0$ is being assumed.

The character $\beta_p$ is nontrivial.  Indeed, if all generating
phases $\lambda_j$ had order prime to $p$, then all matrices of
$J$ and all its character values would belong to
\[
                    E=F(\lambda_1,\ldots,\lambda_r),
\]
a cyclotomic field, or a subfield of a cyclotomic field, with
conductor prime to $p$.  Such a field is unramified at $p$.
For an elementary justification, when $p\nmid M$ the polynomial
$X^M-1$ and its derivative have no common root modulo $p$;
the cyclotomic factors therefore define etale orders locally at
$p$, and their number fields are unramified there.  Subfields
retain that property.  Here $F$ is either $\mathbf Q$ or
$\mathbf Q(\zeta_3)$, so it introduces no ramification at $p$.
But (BI.39) would put $K$ inside $E$, whereas
$\mathbf Q_p(\sqrt{-p})/\mathbf Q_p$ is ramified quadratic by
the Eisenstein polynomial $X^2+p$.  This is impossible.

Because $g_1\cdots g_r=1$, at least two of the $g_j$ have
$\beta_p(g_j)\ne1$.  A nontrivial $p$-primary phase implies
$\mathbf Q(\zeta_p)\subset\mathbf Q(\lambda_j)$, and
\[
 [F(\lambda_j):F]\ge[F(\zeta_p):F]=p-1.
                                                               \tag{BI.44}
\]
For $F=\mathbf Q$ the equality is the usual cyclotomic degree.
For $F=\mathbf Q(\zeta_3)$, the intersection with
$\mathbf Q(\zeta_p)$ is $\mathbf Q$: otherwise the quadratic
field $F$, ramified at three, would lie in a field unramified
at three.  Thus the equality holds in this case as well.
If $g_j$ fixes a subspace of dimension $b$, then
$\lambda_j^{-1}$ is an eigenvalue of $n_j$ of multiplicity $b$.
Its irreducible polynomial over $F$, of degree at least $p-1$,
therefore occurs at least $b$ times in the degree-$e$
characteristic polynomial of $n_j$.  Consequently
\[
 \dim\ker(g_j-1)\le\left\lfloor\frac{e}{p-1}\right\rfloor,
 \qquad
 \operatorname{rank}(1-g_j)\ge
              e-\left\lfloor\frac{e}{p-1}\right\rfloor.
\]
For $p\ge11$ this is full moving rank.  For $p=7$ and $e\le8$
it is at least $e-1$.

The two distinguished inertias have the preceding lower bound,
and every other genuine inertia contributes at least one.  Hence
\[
 2e+6\ge2e-2\left\lfloor\frac{e}{p-1}\right\rfloor+(r-2)
                   \ge2e+8,
\]
a contradiction.  It rules out the original assumption that any
inducing image contains a reflection, and proves the theorem.
\end{proof}
